\section*{Conclusiones}
\addcontentsline{toc}{section}{Conclusiones}

A partir del análisis integral desarrollado a lo largo de la presente monografía sobre la Red Dorsal Nacional de Fibra Óptica (RDNFO), su marco regulatorio, arquitectura técnica, crisis contractual y su impacto socioeconómico en el Perú, se presentan las siguientes conclusiones:

\subsection*{Solidez de la arquitectura técnica frente a la rigidez del modelo contractual}
El diseño técnico e infraestructura física de la RDNFO, basados en una topología de anillos concéntricos, tecnología de multiplexación DWDM y un tendido de 13,571 kilómetros de fibra óptica monomodo, cumplieron con el objetivo de dotar al Perú de una autopista troncal de alta velocidad orientada a conectar las 180 capitales de provincia. Sin embargo, esta solidez tecnológica se vio socavada por un esquema de negocio altamente rígido. El modelo de «Operador de Operadores» sujeto a una tarifa única regulada de US\$ 23 por Mbps (sin IGV) impidió adaptar los precios a la dinámica del mercado, donde la inversión privada independiente redujo drásticamente las tarifas, provocando una subutilización estructural de la red estatal, alcanzando apenas un 8\% de su capacidad operativa.

\subsection*{Desconexión entre el despliegue de infraestructura y la inclusión digital efectiva}
El cumplimiento de las metas de cobertura física no se tradujo automáticamente en el cierre de la brecha digital. Si bien el Perú ha experimentado un incremento notable en las velocidades medianas de descarga de Internet fijo, superando los 200 Mbps hacia 2025, este avance responde fundamentalmente a las inversiones de los operadores privados en sus propias redes de última milla y fibra al hogar (FTTH). Por el contrario, la penetración de Internet en el ámbito rural se mantiene en niveles críticos, alrededor del 23.6\% frente a un 80.5\% en Lima Metropolitana, demostrando que la disponibilidad de una red troncal es insuficiente si no se articula con políticas de incentivo a la demanda y despliegue capilar en zonas de menor rentabilidad.

\subsection*{Inviabilidad del modelo financiero original y dilema de la gestión estatal provisional}
La incapacidad de renegociar las adendas tarifarias con el concesionario Azteca Comunicaciones condujo a la caducidad del contrato por razones de interés público en 2021 y a la reversión de los bienes al Estado. La asunción provisional de la red por parte del Programa Nacional de Telecomunicaciones (PRONATEL) puso en evidencia las limitaciones burocráticas del sector público para administrar un activo tecnológico de tal magnitud, generando un elevado déficit fiscal con costos operativos cercanos a los US\$ 70 millones frente a recaudaciones reducidas y acelerando el riesgo de obsolescencia tecnológica de los nodos. El esquema de «Pago Cero» ha servido como un paliativo social para conectar instituciones públicas, pero no resuelve la sostenibilidad financiera ni comercial de la red a largo plazo.

\subsection*{Subutilización de servicios públicos estratégicos e impacto social postergado}
Los fines sociales de alto impacto proyectados en la Ley N.° 29904 se han concretado de manera marginal y fragmentada. Iniciativas clave como la Red Nacional de Investigación y Educación (RNIE) han sufrido severos retrasos que limitaron la integración académica e internacional de las universidades del interior del país, mientras que las experiencias en telesalud y gobierno digital se mantuvieron mayoritariamente como proyectos piloto. Este desaprovechamiento representó un elevado costo de oportunidad para el desarrollo humano de las poblaciones más vulnerables del Perú profundo.

\subsection*{Reorientación estratégica: Flexibilidad, alianzas y priorización de la última milla}
Las lecciones aprendidas del caso peruano y las tendencias comparadas en América Latina evidencian que el futuro de la RDNFO no requiere expandir indefinidamente la red de transporte estatal, sino replantear su modelo de gestión. La optimización del proyecto demanda la introducción urgente de un régimen de flexibilización tarifaria que contemple descuentos por volumen y diferenciación geográfica, la estructuración de nuevas Alianzas Público-Privadas (APP) bajo modelos de gobernanza adaptables, y una articulación eficiente con mecanismos regulatorios de última milla, como el Canon por Cobertura, para lograr que la infraestructura troncal ya desplegada cumpla su propósito fundacional: democratizar el acceso a la banda ancha en todo el territorio nacional.